\chapter{相关研究文献综述}

由于目前的相关文献对于教育模式对大学生创业意愿的影响研究还存在一定的空白，本研究选取了天津市某高校创业学院选修创新创业课程的学生作为调查对象，探究教育模式在培育大学生创业意识方面的作用机理。研究把教育模式分为理论型、实践型和整合型三类，把自主动机当作核心中介变量，又纳入了外部激励因素作为调节变量，创建起了\nolinebreak[4]“教育模式→自主动机→创业意愿”以及\nolinebreak[4]“外部激励×自主动机→创业意愿”的两条路径模型，探究外部激励怎样经由调节自主动机来影响人的创业行为倾向。以实证分析的形式，探究不同种类的教育教学方法怎样借助不同的动机激发机制来影响大学生的创业意向，给高校创新创业人才的培养体系改进赋予实证支撑和操作参照。

\section{创新创业教育模式}

\subsection{创新创业教育的内涵与发展}

创新创业教育是以培养具备创新精神和实践能力的高素质人才为根本目标的系统性育人模式。王洪才\nolinebreak[4]（2021）指出，这是中国特色高等教育体系的重要组成部分，其本质在于突破传统应试教育的束缚，着重培养学生的主体意识和创造能力\cite{wanghongcai2021}。与传统职业教育侧重于职业技能培养不同，创新创业教育更加注重学生在机会识别、资源整合、风险控制等方面的综合素养形成\cite{zhoudongmei2020}。

在国际化背景下，创新创业教育已成为高等教育系统变革的重要驱动力。关于创业教育在高等教育中的影响，Nabi等\nolinebreak[4]（2017）的系统综述指出，创业教育项目的效果受到项目设计、实施方式和学生个体特征的多重影响\cite{Nabi2017}。对创业教育与学生创业机会识别能力的系统文献综述\cite{EEP2025a}以及对加纳大学学生创业者的定性研究\cite{EEP2025b}都表明，创业教育的效果需要在具体情境中加以评估。邴浩\nolinebreak[4]（2019）的实证研究表明，创业教育的效果与个体特征密切相关，并非对所有学生都能产生相同的促进作用\cite{binghao2019}。

近年来，产教融合成为提升创新创业人才素质的主要路径。基于社会全面参与的创业教育模式研究\cite{guanlixuebao2011}发现，教育模式的创新需要充分考虑社会各界的参与和支持。大学生创业者心理所有权对警觉性市场学习的影响\cite{guanlixuebao2023}以及\nolinebreak[4]“首届创业学与企业家精神教育研讨会”会议综述\cite{nankai2003}的发现，进一步强调了教育模式差异化设计的重要性。

\subsection{创新创业教育模式的类型划分}

依据教学目标和实施路径的不同，创新创业教育模式呈现出多样化特征。Manimala和Thomas\nolinebreak[4]（2017）将创业教育分为三类：理论型、实践型以及两者结合的融合型\cite{Manimala2017}。

在创业教育的具体实践形式方面，Gielnik等\nolinebreak[4]（2015）评估了学生创业培训的效果，发现行动和行动调节在创业过程中发挥着关键作用\cite{Gielnik2015}。Gielnik等\nolinebreak[4]（2017）从长期视角探讨了创业培训的效果，发现培训能够持续激发学生的创业激情\cite{Gielnik2017}。然而，Bohlayer和Gielnik\nolinebreak[4]（2023）也指出，行动导向的创业培训可能导致创业自我效能感的暂时性下降，这被称为\nolinebreak[4]“训练困境”\cite{Bohlayer2023}。

国内高校在推行创新创业教育过程中呈现出较为显著的类型化特征。本研究以天津市某高校创业学院课程设置为案例，将创新创业教育实践形态归纳为三种主要模式：理论型课程主要采用课堂讲授式教学方法，侧重于创业学的理论体系、运作机制和模式建构等内容传授；实践型课程通常采用创业实训、项目驱动和仿真竞赛等方式，使学生置身于真实或者模拟的环境中进行系统的创业能力训练；融合型课程是理论教学和实践教学相结合的一种新的教学方式，体现着系统性的理论讲授同综合性的实践活动相融合的特点。

\subsection{创新创业教育的效果研究}

学界对创新创业教育的效果已积累丰富研究成果。Souitaris等\nolinebreak[4]（2007）采用准实验研究方法发现，创业教育通过知识传授、思维启迪、资源支持三条路径促进理工科学生的创业意愿形成\cite{Souitaris2007}。Rippa等\nolinebreak[4]（2020）的量化分析表明，创业课程强度虽非影响创业意愿的主要因素，但对工程类学生创业意愿的形成具有重要作用\cite{Rippa2020}。Nowiński等\nolinebreak[4]（2017）基于波兰维谢格拉德地区的跨区域比较研究发现，创新创业教育能够提升学习者的自我效能感，进而增强创业动机\cite{Nowinski2017}。

国内研究方面，罗佳等\nolinebreak[4]（2025）通过实证研究发现，创业教育能够提升大学生的数字创业意愿，但社交媒体的使用具有负向调节作用\cite{luojia2025}。李婧和郭佳玮\nolinebreak[4]（2025）通过准自然实验证实，数据要素集聚为高校学生创新创业活动提供了技术支撑和资源保障\cite{lijing2025}。郑馨等\nolinebreak[4]（2025）运用跨国组态分析框架发现，政府干预、市场机制和社会制度三者之间的互动配合是激发高校学生高潜力创业行为的关键因素\cite{zhengxin2025}。李其容等\nolinebreak[4]（2023）发现，创业自我效能在创业进展与创业努力之间发挥关键中介作用，且这一中介效应受个体调节定向的调节\cite{liqirong2023}。

尽管学界已普遍认同创业教育的综合性特征，但对不同课程模式功能特点及内在差异的深入研究仍显不足。这一学术空白既制约了具体实践研究的深入开展，也影响了理论指导作用的充分发挥。

\section{创业意愿}

\subsection{创业意愿的概念界定}

创业意愿指的是个人对未来自己是否自主创业或者参加创业活动有主观倾向以及实现的可能性\nolinebreak[4]（Liñán \& Chen，2009）\cite{Linan2009}。按照态度和行为互相影响的理论架构，此变量经常被看作是预示个体后续创业行为的决定性前期指标。

胡望斌等\nolinebreak[4]（2025）在对创业型领导对员工内创业行为影响的研究中发现，创业意愿是连接领导行为与员工创业行为的关键心理机制\cite{huwangbin2025}。郭润萍等\nolinebreak[4]（2024）通过对数字创意新企业的多案例研究，揭示了社会互动视角下数字创业机会的客观化机理，其中创业意愿发挥着重要的驱动作用\cite{guorunping2024}。

\subsection{创业意愿的测量}

创业意愿测评工具经过长期发展，已经形成较为完善的体系。Liñán和Chen\nolinebreak[4]（2009）根据计划行为理论创建的\nolinebreak[4]《创业意向问卷》\nolinebreak[4]（EIQ），现在已经成为学术界广泛应用的经典量表之一\cite{Linan2009}。经过跨文化多国实证研究显示，本量表在世界各国数据检验中具有很好的信度与效度，涵盖意图强度、目标导向性、准备程度、风险感知等主要方面，可以用来对个人的创业倾向进行系统性的评价。

在创业意向的测量方面，Al-Jubari等\nolinebreak[4]（2019）整合自我决定理论与计划行为理论，使用了经过验证的测量工具对创业意向进行评估\cite{AlJubari2019}。Kazmi和Nábrádi\nolinebreak[4]（2017）在研究中同样采用标准化的量表来测量学生的创业意向和行为变化\cite{Kazmi2017}。

\subsection{创业意愿的影响因素}

创业意愿形成是个体因素与环境因素交互作用的结果。在个体因素方面，姜诗尧等\nolinebreak[4]（2024）从生命史理论视角出发，探讨了创业者童年社会经济地位对突破式创新的影响机制\cite{jiangshiyao2024}。叶文平等\nolinebreak[4]（2017）发现，财富积累速度与创业者进取心之间存在显著关联，制度环境感知在其中发挥调节作用\cite{yewenping2017}。

在环境因素方面，Edelman等\nolinebreak[4]（2016）的研究揭示了家庭支持对年轻创业者创业活动的重要影响\cite{Edelman2016}。Kaandorp等\nolinebreak[4]（2020）对学生创业者的初始网络过程研究发现，社会网络的构建对学生创业意愿的形成具有重要作用\cite{Kaandorp2020}。胡玲玉等\nolinebreak[4]（2014）研究探讨了创业环境和创业自我效能对个体创业意向的影响，发现环境因素通过自我效能感间接影响创业意向\cite{hulingyu2014}。高校毕业生创业意向影响因素研究发现，不同类型的激励方式对学生创业意愿的影响存在显著差异\cite{wangchunwen2024}。

Al-Jubari等\nolinebreak[4]（2019）整合自我决定理论与计划行为理论，发现自主、胜任和归属需求通过态度、主观规范和知觉行为控制间接影响创业意向，模型解释了71\%的方差\cite{AlJubari2019}。Kazmi和Nábrádi\nolinebreak[4]（2017）提出，创业教育通过创业自我效能和感知合意性影响意向，\nolinebreak[4]“为创业而教”比\nolinebreak[4]“关于创业而教”效果更强\cite{Kazmi2017}。

\section{自主动机与外部动机}

\subsection{自我决定理论与动机分类}

自我决定理论最早是由Deci和Ryan在1985年提出的一个学术框架，它主要研究人的行为动机的本质以及这些动机对于人的心理发展所起的作用。把动机分成自主型和管控型这两种基本类型，自主型动机来源于人的内在兴趣倾向、价值取向、全面的自我认同建构，在具体的行动中体现出明显的主观能动性、独立选择性；管控型动机受外界压力、强化激励或者社会规范所制约，一般表现出来的是顺应性的行为模式。

在此基础上，Vallerand等人\nolinebreak[4]（1992）创建了学术动机量表\nolinebreak[4]（AMS），将其分成内部动机、认同调节、内摄调节、外部调节和无动机五个连续的学术动机维度\cite{Vallerand1992}，如表\ref{tab:ams_dim}所示。其中，内部动机和认同调节属于自主型动机；内摄调节、外在调节和无动机都被视为受控型动机。

\tablemacro{htbp}
            {学术动机维度表}
            {ccp{6cm}}
            {\songtibf{动机类型} & \songtibf{自主程度} & \songtibf{说明} \\}
            {内部动机 & 最高 & 因活动本身的乐趣和满足感而参与 \\
            认同调节 & 较高 & 因认同活动的价值和意义而参与 \\
            内摄调节 & 中等 & 因内在压力\nolinebreak[4]（如内疚、焦虑）而参与 \\
            外部调节 & 较低 & 因外部奖惩\nolinebreak[4]（如学分、奖励）而参与 \\
            无动机 & 最低 & 缺乏参与的理由或意图 \\}
            {tab:ams_dim}

这五种动机类型可以构成一个由高自主性到高控制性的连续谱系模型。其中，内部动机和认同调节属于自主型动机；内摄调节、外在调节和无动机都被视为受控型动机，即本研究所用的外部型动机。根据自我决定理论的实证研究发现，胜任需求得到满足以后，可以经由加强自我效能感来推动创业动机的增长，还可以明显提升学习行为的积极性。与此同时，Al-Jubari等\nolinebreak[4]（2019）整合自我决定理论与计划行为理论，验证了自主需求、胜任需求和归属需求通过态度、主观规范和知觉行为控制间接影响创业意向的路径\cite{AlJubari2019}。这一研究为本研究探讨自主动机的中介作用提供了坚实的理论基础。

\subsection{自主动机与外部动机的测量}

运用自我决定理论框架的量化评价工具，对自主动机、外在动机表现形式做系统的考察。学术动机量表\nolinebreak[4]（AMS）是Vallerand等人在1992年开发的，它被广泛使用，包含五个方面即内在动机、认同调节、内化调节、外在调节和无动机\cite{Vallerand1992}。根据研究目的，本研究从AMS精心挑选具有代表性的题目来准确地测量出这五种动机类型所表现出的具体行为。

自主动机是内在驱动力和认同调节机制构成的，有五个主要方面，即兴趣倾向、心理满足感、自我突破与价值实现具体目标、创业技能对未来职业生涯发展的重要意义、提高个人综合素养。

外部动机分为内摄调节、外在规劝、无动机这三类，共5个量表题目。M6用来评价个人在课程中的内在需求参与度，M7考察的是学生以展示为课程展示方式的具体行为，M8是学生成果为学分的动机导向的表达情况，M9反映的是为了得到奖励、取得资格认证而产生的学习动机，M10是部分学生对于课程目标认识不清或者缺乏学习动机的状态。这些指标一起形成了对外部动机类型系统的测量体系。

本研究把自主动机当作主要的中介变量，创建起创新创业教育模式、自主动机和创业意愿三者间的理论模型。从多种类型的课程入手，探究其对学生的内生学习兴趣以及价值认同产生怎样的影响，并从教育实践角度出发，探究个体创业意愿形成的一些规律。引入外部动机做调节变量，考察它的边界效应在创业意愿形成的全过程里，把认知行为理论加进去，加深对激励约束机制作用路径和强弱差别方面的认识。

\subsection{自主动机、外部动机与创业意愿的关系}

本研究主要研究自主动机同创业意愿之间的联系。按照自我决定理论的分析框架，当个人参与创业相关的学习活动时，自主性倾向越强，认知投入就越深入，情感投入就越强烈，行为执行就越高效，创业意愿也就越可能得到培养和增强。相反的，如果一个人具有受控导向的行为特征，那么就会影响人们对创业价值的认识程度以及建立自我观念的过程。

在国际研究方面，Al-Jubari等\nolinebreak[4]（2019）的实证研究表明，自主动机对创业意向具有显著的预测作用，这一效应通过态度、主观规范和知觉行为控制的中介作用实现\cite{AlJubari2019}。Kazmi和Nábrádi\nolinebreak[4]（2017）的研究发现，创业教育的效果受到学生动机类型的调节，\nolinebreak[4]“为创业而教”比\nolinebreak[4]“关于创业而教”更能激发学生的自主动机\cite{Kazmi2017}。

关于创业激情的研究也为理解自主动机的作用提供了重要启示。李其容等\nolinebreak[4]（2024）基于身份控制理论，探究了创业激情的两个组成部分——创业身份中心性和强烈积极情绪——的不同匹配状态对创业努力和创业成瘾的影响\cite{liqirong2024}。研究发现，二者匹配时创业努力更高、成瘾更低；而\nolinebreak[4]“错位”状态则通过自满感和焦虑感等成就情绪产生影响。这一发现揭示了情感因素在创业动机形成过程中的重要作用。

\section{研究述评}

综合上述文献分析，当前创新创业教育研究存在以下不足，构成本研究的主要切入点。

第一，创新创业教育模式研究多集中于宏观层面，对不同类型教育模式差异化特征的系统研究较为薄弱。现有文献大多将创业教育视为单一变量，关注其对创业意愿的整体影响，较少关注不同教育方法的价值实现路径和适用场景差异。尽管Manimala和Thomas\nolinebreak[4]（2017）已构建初步的教育模式分类框架\cite{Manimala2017}，但后续实证研究多采用\nolinebreak[4]“参与与否”或\nolinebreak[4]“参与程度”等简单测量指标，未能充分展现各类教育方式的独特性及其影响机理。

第二，自主动机作为中介变量的研究尚处于起步阶段。尽管已有研究者将自主动机纳入分析框架，但对其具体传导机制的研究尚不充分。Al-Jubari等\nolinebreak[4]（2019）的研究虽然证实了自主动机对创业意向的影响\cite{AlJubari2019}，但未深入探讨教育模式如何通过自主动机影响创业意愿的中介路径。不同教育模式下自主动机中介效应的异质性比较研究仍较为缺乏。

第三，外部动机调节作用的研究尚未形成系统的本土化体系。在创业教育情境下，外部激励因素如何影响学习者的自主动机与创业意愿之间的互动关系，尚缺乏系统且有实证支撑的论证。现有文献多探讨外部动机的线性效应或独立作用路径，对其作为潜在调节变量的作用机理探讨不足。

综上所述，本研究选取天津市某高校创业学院选修创新创业课程的学生为调查对象，将教育模式划分为理论型、实践型和融合型三类，以自主动机为核心中介变量，纳入外部动机为调节变量，构建\nolinebreak[4]“教育模式→自主动机→创业意愿”及\nolinebreak[4]“外部动机×自主动机→创业意愿”的双路径模型，通过实证分析探究不同教育教学方法如何经由不同的动机激发机制影响大学生的创业意向，为高校创新创业人才培养体系的优化提供实证支撑与实践参照。

\clearpage
